% ===== tr-05 块1b：第 4 章 pp.32--35（原文页 32--35）=====
% [原文 p.32]
上面的陈述只是让人对实际的真相有个感觉。证明以及精确得多的陈述最终将出现在 [30] 中。

若 $\mathcal{X}\to D$ 是\emph{极化的}环面退化，即存在 $\mathcal{X}$ 上的相对丰沛线丛 $\mathcal{L}$，则我们可以构造另一个整热带流形 $(\check B,\check{\mathcal{P}})$，称之为\emph{相交复形}，如下。

对 $\mathcal{X}_0$ 的每个不可约分支 $X_i$，$\mathcal{L}|_{X_i}$ 是环面簇上的丰沛线丛。令 $\check\sigma_i\subseteq N_{\mathbb{R}}$ 记该线丛的 Newton 多胞形。于是，含于 $X_i$ 中的 $\mathcal{X}_0$ 的分层与 $\check\sigma_i$ 的面之间存在一一的保持包含的对应。我们便可以按显然的方式把诸 $\check\sigma_i$ 粘合起来：若 $Y$ 是 $\mathcal{X}_0$ 的余维一分层，则它含于两个不可约分支 $X_i$ 与 $X_j$ 之中，并定义 $\check\sigma_i$ 与 $\check\sigma_j$ 的面。这些面仿射同构，因为它们都是 $\mathcal{L}|_Y$ 的 Newton 多胞形，于是我们可以按标准方式把它们等同起来。这样便得到一个带多面体分解 $\check{\mathcal{P}}$ 的拓扑空间 $\check B$。

要定义顶点 $v\in\check{\mathcal{P}}$ 处的扇结构，注意这样的顶点对应于 $\mathcal{X}_0$ 的一个零维分层，从而给出对偶相交复形的一个极大胞腔 $\sigma_v$。取 $v$ 处的扇结构为用 $\sigma_v$ 的法扇 $\check\Sigma_v$ 定义的结构。于是，$\check\Sigma_v$ 中的锥与含有 $v$ 所对应分层的 $\mathcal{X}_0$ 的分层之间存在一一的保持包含的对应。这一对应使我们能定义扇结构
\[
S_v:U_v\to\mathbb{R}^n,
\]
它把 $U_v\cap\check\sigma$（对任意以 $v$ 为顶点、属于 $\check{\mathcal{P}}$ 的 $\check\sigma$）映入 $\check\Sigma_v$ 中相应的锥。容易验证这组扇结构满足整热带流形的定义，从而定义了相交复形 $(\check B,\check{\mathcal{P}})$。

与定理 4.7.8 类似，我们期待

\begin{conjecture}{4.7.9}
设 $\mathcal{X}\to D$ 是极化环面退化，相交复形为 $(\check B,\check{\mathcal{P}})$。设 $\omega_t$ 是 $X_t$ 上代表该极化的第一 Chern 类的 Kähler 形式。则存在开集 $\check U\subseteq\check B$，使得 $\check B\setminus\check U$ 收缩到 $\check B$ 的判别轨迹 $\Gamma$ 上，并且对任意 $t$，$X_t$ 都是 $\check X(\check U)$ 的辛紧化。
\end{conjecture}

我并不认为这特别困难：它应当可以用 W.-D. Ruan [71] 的技术来处理，但这样的进路在一般情形下尚未被实现。

相交复形与对偶相交复形之间的关系，可以通过引入多值分段线性函数而变得更精确，这与定义 4.1.3 中的多值凸函数类似。

\begin{definition}{4.7.10}
设 $(B,\mathcal{P})$ 是整热带流形。则 $B$ 上的一个\emph{多值分段线性函数} $\varphi$ 是开覆盖 $\{(U_i,\varphi_i)\}$ 上的一族连续函数，使得 $\varphi_i$ 在每个与 $U_i$ 相交的 $\mathcal{P}$ 的胞腔上仿射线性，且在 $U_i\cap U_j$ 上，
% [原文 p.33]
% [Tutor Note: 「整斜率」条件原文写作「differential in $\check\Lambda$」，用的是带 check 的对偶格 $\check\Lambda$，不是 $\Lambda$；抄写时勿丢掉 check。]
$\varphi_i-\varphi_j$ 是仿射线性的。此外，对任意 $\tau\in\mathcal{P}$，设 $S_\tau:U_\tau\to\mathbb{R}^k$ 是诱导出的扇结构。则扇 $\Sigma_\tau$ 上存在一个分段线性函数 $\varphi_\tau$，使得在 $U_i\cap U_\tau$ 上，$\varphi_i-\varphi_\tau\circ S_\tau$ 是仿射线性的。这里我们将始终假定每个 $\varphi_i$ 的线性部分的微分属于 $\check\Lambda$，即 $\varphi_i$ 具有整斜率。
\end{definition}

关于每个 $\varphi_i$ 在 $U_\tau$ 上局部性态的这个颇为技术性的条件，来自如下想法：这样的多值分段线性函数其实不过是由 $(B,\mathcal{P})$ 的扇结构所给出的诸扇 $\Sigma_\tau$ 上一族分段线性函数；这些函数需要满足某些相容性条件，而这种相容性是由下面的讨论所驱动的。

假设给定一个极化环面退化 $\mathcal{X}\to D$。事实上，我们如下得到对偶相交复形 $(B,\mathcal{P})$ 上的一个多值分段线性函数 $\varphi$。限制到任意环面分层 $X_\tau$，$\mathcal{L}|_{X_\tau}$ 完全由 $\Sigma_\tau$ 上的一个整分段线性函数 $\varphi_\tau$ 确定，且在相差一个线性函数的意义下是唯一确定的。把这一分段线性函数经由 $S_\tau$ 拉回到 $U_\tau$，我们得到一族分段线性函数 $\{(U_\tau,\varphi_\tau\circ S_\tau)\mid\tau\in\mathcal{P}\}$。由 $\tau\subseteq\sigma$ 时 $(\mathcal{L}|_{X_\tau})|_{X_\sigma}=\mathcal{L}|_{X_\sigma}$ 可知，在重叠区域上 $\varphi_\sigma\circ S_\sigma$ 与 $\varphi_\tau\circ S_\tau$ 至多相差一个仿射线性函数。于是 $\{(U_\tau,\varphi_\tau\circ S_\tau)\}$ 定义了一个多值分段线性函数。多值分段线性函数定义中的最后一个条件，则反映了如下需要：在 $\sigma$ 的邻域内，该函数局部地是经由 $S_\sigma$ 拉回的某个函数。

若 $\mathcal{L}$ 丰沛，则由 $\mathcal{L}|_{X_\sigma}$ 确定的分段线性函数是严格凸的。因此我们说一个多值分段线性函数是\emph{严格凸}的，如果对每个 $\tau\in\mathcal{P}$，$\varphi_\tau$ 都严格凸。

于是，若 $\mathcal{X}\to D$ 是极化环面退化，我们将用 $(B,\mathcal{P},\varphi)$ 表示对偶相交复形以及诱导出的多值函数 $\varphi$ 这一组数据。我们称这个三元组为该极化退化的对偶相交复形。

现在假设抽象地给定三元组 $(B,\mathcal{P},\varphi)$，其中 $(B,\mathcal{P})$ 是整热带流形，$\varphi$ 是 $B$ 上的严格凸多值分段线性函数。然后我们如下构造 $(B,\mathcal{P},\varphi)$ 的\emph{离散 Legendre 变换} $(\check B,\check{\mathcal{P}},\check\varphi)$。

$\check B$ 将通过粘合 Newton 多胞形来构造。若对 $\mathcal{P}$ 的顶点 $v$，把扇 $\Sigma_v$ 看作生活在 $M_{\mathbb{R}}$ 中，则 $\varphi_v$ 的 Newton 多胞形为
\[
\check v:=\{x\in N_{\mathbb{R}}\mid \langle x,y\rangle\ge-\varphi_v(y)\ \ \forall y\in M_{\mathbb{R}}\}.
\]
$\check v$ 的面与含 $v$ 的 $\mathcal{P}$ 的胞腔之间存在一一的包含反序对应。此外，若 $\sigma$ 是 $\mathcal{P}$ 中含两个顶点 $v$ 与 $v'$ 的最小胞腔，则 $\check v$ 与 $\check v'$ 的相应面整仿射同构，因为它们都同构于 $\varphi_\sigma$ 的 Newton 多胞形。于是我们可以沿这个公共面把 $\check v$ 与 $\check v'$ 粘合起来。作完所有这些等同之后，我们得到一个胞腔复形 $(\check B,\check{\mathcal{P}})$，它
% [原文 p.34]
其实正是 $(B,\mathcal{P})$ 的对偶胞腔复形。它的整热带结构由如下方式给出：对 $\sigma\in\mathcal{P}_{\max}$，取顶点 $\check\sigma$ 处的扇结构为 $\sigma$ 的法扇。

最后，函数 $\varphi$ 在 $(\check B,\check{\mathcal{P}})$ 上有一个离散 Legendre 变换 $\check\varphi$。对于与极大胞腔 $\sigma\in\mathcal{P}$ 对偶的顶点 $\check\sigma\in\check{\mathcal{P}}$，我们别无选择，只能把 $\check\varphi$ 在其邻域内定义为以 $\sigma$ 为 Newton 多胞形的分段线性函数，即
% [Tutor Note: 定义 $\check\varphi_{\check\sigma}(y)=-\inf\{\langle y,x\rangle\mid x\in\sigma\}$ 中的负号不可漏；因 $\sigma$ 紧，$\inf$ 实际可取为 $\min$，但原文写的是 $\inf$，正文照抄。]
\[
\check\varphi_{\check\sigma}(y)=-\inf\{\langle y,x\rangle\mid x\in\sigma\subseteq M_{\mathbb{R}}\}.
\]
这样便得到 $(\check B,\check{\mathcal{P}},\check\varphi)$，即 $(B,\mathcal{P},\varphi)$ 的离散 Legendre 变换。若 $B$ 是 $\mathbb{R}^n$，则它与经典的离散 Legendre 变换概念一致。离散 Legendre 变换具有若干有用的性质：
\begin{itemize}
\item $(\check B,\check{\mathcal{P}},\check\varphi)$ 的离散 Legendre 变换是 $(B,\mathcal{P},\varphi)$。
\item 若我们把底拓扑空间 $B$ 与 $\check B$ 通过对偶胞腔复形的底空间而等同起来，则 $\Lambda_{B_0}\cong\check\Lambda_{\check B_0}$ 且 $\check\Lambda_{B_0}\cong\Lambda_{\check B_0}$，其中下标指明是用哪一个仿射结构来定义 $\Lambda$ 或 $\check\Lambda$ 的。
\end{itemize}

这有望清楚地表明：离散 Legendre 变换是第 2 节中由 Legendre 变换所提供的对偶性的合适替代。

特别地注意，若 $\mathcal{X}\to D$ 是极化环面退化，对偶相交复形为 $(B,\mathcal{P},\varphi)$，则其离散 Legendre 变换 $(\check B,\check{\mathcal{P}},\check\varphi)$ 满足：$(\check B,\check{\mathcal{P}})$ 正是该极化退化的相交复形。函数 $\check\varphi$ 是 $\check B$ 上的一些额外信息，由离散 Legendre 变换的定义，它编码了 $\mathcal{P}$ 的胞腔。对偶相交复形的这些胞腔是用 $\mathcal{X}\to D$ 的局部环面结构定义的。因此 $\check\varphi$ 可以看作携带了关于这种局部环面结构的信息。我们将称 $(\check B,\check{\mathcal{P}},\check\varphi)$ 为极化环面退化 $\mathcal{X}\to D$ 的相交复形。

于是我们看到：对 $(B,\mathcal{P},\varphi)$ 而言，$\mathcal{P}$ 携带关于对数结构的信息，而 $\varphi$ 携带关于极化的信息；但对 $(\check B,\check{\mathcal{P}},\check\varphi)$ 而言，$\check{\mathcal{P}}$ 携带关于极化的信息，而 $\check\varphi$ 携带关于对数结构的信息。镜像对称把这两类信息互换！

现在我们可以陈述一个\emph{代数几何的 SYZ 程序}。仿照第 5 节所建议的程序，我们可以遵循以下步骤：
\begin{enumerate}
\item 从带有丰沛极化的 Calabi--Yau 流形环面退化 $\mathcal{X}\to D$ 出发。
\item 由这些数据构造对偶相交复形 $(B,\mathcal{P},\varphi)$，如上所述。
\item 作离散 Legendre 变换，得到 $(\check B,\check{\mathcal{P}},\check\varphi)$。
\item 尝试构造 Calabi--Yau 流形的极化退化 $\check{\mathcal{X}}\to D$，使其对偶相交复形为 $(\check B,\check{\mathcal{P}},\check\varphi)$，或其相交复形为 $(B,\mathcal{P},\varphi)$。
\end{enumerate}
% [原文 p.35]
% [Tutor Note: 例 4.7.11 中两个求和分别取在 $m\in M\cap\Delta$ 与 $n\in N\cap\nabla$（格 $M,N$，不是 $M_{\mathbb R},N_{\mathbb R}$），抄写时勿写成 $M_{\mathbb R}\cap\Delta$；裁图 p4 p.35 上方可核对。]
\begin{example}{4.7.11}
离散 Legendre 变换使我们能够重现 Batyrev 对偶 [5]。设 $\Delta\subseteq M_{\mathbb{R}}$ 是自反多胞形，$\nabla\subseteq N_{\mathbb{R}}$ 是它的极对偶，并设 $0\in\Delta$ 是唯一的内点。于是我们得到两个环面退化，分别由方程
\[
s_0+t\sum_{m\in M\cap\Delta}c_ms_m=0,\qquad
s_0+t\sum_{n\in N\cap\nabla}c_ns_n=0
\]
在 $\mathbb{P}_\Delta\times\mathbb{A}^1$ 与 $\mathbb{P}_\nabla\times\mathbb{A}^1$ 中给出，其中 $s_m$ 是 $\mathcal{O}_{\mathbb{P}_\Delta}(1)$ 中对应于 $m$ 的截面（$s_n$ 是 $\mathcal{O}_{\mathbb{P}_\nabla}(1)$ 中对应于 $n$ 的截面）。容易验证这两个退化的对偶相交复形如下：对第一个退化，$B=\partial\nabla$，多面体分解由 $\nabla$ 的真面给出，每个顶点 $v$ 处的扇结构由投影 $U_v\hookrightarrow N_{\mathbb{R}}\to N_{\mathbb{R}}/\mathbb{R}v$ 给出；对第二个退化，则改用 $\Delta$ 代替 $\nabla$。然后可以验证：若分别用 $\mathcal{O}_{\mathbb{P}_\Delta}(1)$ 与 $\mathcal{O}_{\mathbb{P}_\nabla}(1)$ 极化这两个退化，则相应的三元组 $(B,\mathcal{P},\varphi)$ 互为 Legendre 对偶。于是 Batyrev 对偶是这一镜像构造一般进路的特例。
\end{example}

对于更一般的构造------它适用于环面簇中完全交 Calabi--Yau 的镜像的 Batyrev--Borisov 构造 [6]------见 [22]。

这一镜像对称算法中唯一缺少的步骤是最后一步：

\begin{question}{4.7.12}[重构问题，版本 II]
给定 $(B,\mathcal{P},\varphi)$，能否构造一个极化环面退化 $\mathcal{X}\to D$，使其相交复形为 $(B,\mathcal{P},\varphi)$？
\end{question}

人们或许希望通过朴素的形变理论来解决这个问题：由数据 $(B,\mathcal{P},\varphi)$ 构造中心纤维 $\mathcal{X}_0$，然后形变它以求得到光滑化。然而，正如 Kawamata 与 Namikawa 最初在正规交叉情形所观察到的 [51]，必须先给 $\mathcal{X}_0$ 附加某种额外结构，它才具有良好的形变理论。这一结构就是\emph{对数结构}；引入对数结构使我们能够直接在退化纤维本身上研究镜像对称的许多方面。因此，让我们转而回顾对数结构理论。

\section{对数结构}
我们回顾 Fontaine--Illusie 与 Kato 的对数结构概念（[45]、[50]）。这些概念在试图通过退化理解镜像对称时起着关键作用。

\begin{definition}{4.8.1}
概形（或解析空间）$X$ 上的一个对数结构，是（乘法、交换）幺半群层的一个（含单位的）同态
\[
\alpha_X:\mathcal{M}_X\to\mathcal{O}_X,
\]
它诱导出同构 $\alpha_X^{-1}(\mathcal{O}_X^\times)\to\mathcal{O}_X^\times$。$\mathcal{O}_X$ 上的幺半群结构由乘法给出。三元组 $(X,\mathcal{M}_X,\alpha_X)$ 于是称为一个\emph{对数空间}。我们常把整个这一套简记为 $X^\dagger$。
\end{definition}
